% Generated draft: requires mathematical and source review.
\documentclass[11pt]{article}
\usepackage[margin=25mm]{geometry}
\usepackage{fontspec}
\setmainfont{Latin Modern Roman}
\usepackage{amsmath,amssymb,mathtools,unicode-math}
\setmathfont{Latin Modern Math}
\usepackage{enumitem,needspace,xurl,hyperref}
\setlength{\parindent}{0pt}
\setlength{\parskip}{5pt}
\setlength{\emergencystretch}{3em}
\newcommand{\N}{\mathbb{N}}
\newcommand{\Z}{\mathbb{Z}}
\newcommand{\Q}{\mathbb{Q}}
\newcommand{\R}{\mathbb{R}}
\newcommand{\C}{\mathbb{C}}
\newcommand{\F}{\mathbb{F}}
\newcommand{\A}{\mathbb{A}}
\newcommand{\PP}{\mathbb{P}}
\newcommand{\E}{\mathbb{E}}
\newcommand{\eps}{\varepsilon}
\newcommand{\dd}{\,\mathrm{d}}
\begin{document}
\section*{UnsolvedMath ID: 30002612}
\textbf{Source problem code:} catalogue entry
\subsection{Definitions and mathematical statement}
The supplied question is the Beauville--Voisin conjecture for projective hyperkähler manifolds. Its coefficient convention and the precise meaning of ``hyperkähler'' are not explicitly specified. The following is a precise formulation under the conventional interpretation of complex irreducible holomorphic symplectic manifolds and rational Chow groups; these conventions require confirmation against the original source.

Let \(X\) be a smooth connected projective complex variety of dimension \(d=2n\), with \(n\geq 1\), which is simply connected and satisfies
\[
H^0(X,\Omega_X^2)=\C\omega
\]
for an everywhere nondegenerate holomorphic two-form \(\omega\). Write
\[
\mathrm{CH}^k(X)_{\Q}=\mathrm{CH}^k(X)\otimes_{\Z}\Q
\]
for codimension-\(k\) algebraic cycles modulo rational equivalence with rational coefficients, and let \(\mathrm{CH}^*(X)_{\Q}\) carry its intersection product. Define the graded subalgebra
\[
R^*(X)=\Q\bigl[c_1(L),\ c_i(T_X)\mid L\in\operatorname{Pic}(X),\ 1\leq i\leq d\bigr]
\subseteq \mathrm{CH}^*(X)_{\Q},
\]
where \(T_X\) is the tangent bundle and \(c_i\) denotes its \(i\)-th Chern class. Thus each element is a polynomial involving finitely many of these generators.

A class \(z\in\mathrm{CH}^k(X)_{\Q}\) is numerically trivial if
\[
\deg(z\cdot\alpha)=0
\qquad\text{for every }\alpha\in\mathrm{CH}^{d-k}(X)_{\Q}.
\]
The question, under the stated conventions, is whether for every such \(X\), every integer \(0\leq k\leq d\), and every \(z\in R^k(X)\),
\[
\left(\forall\alpha\in\mathrm{CH}^{d-k}(X)_{\Q},\quad
\deg(z\cdot\alpha)=0\right)
\quad\Longrightarrow\quad
z=0\text{ in }\mathrm{CH}^k(X)_{\Q}.
\]
For a polynomial with components in several codimensions, numerical triviality and the conclusion are interpreted componentwise. Vanishing here means rational equivalence to zero with rational coefficients. The source does not resolve whether an integral-coefficient assertion, or a broader convention for hyperkähler manifolds, was intended.
\subsection{Short English statement}
For every projective hyperkähler manifold, must a numerically trivial polynomial in divisor classes and tangent-bundle Chern classes be rationally equivalent to zero? The displayed formulation uses complex irreducible holomorphic symplectic manifolds and rational coefficients, subject to confirmation of these unstated source conventions.
\subsection{Sources}
\begin{itemize}
\item Supplied local UnsolvedMath registry record, permanent ID \(30002612\), exact source problem code \texttt{catalogue entry}, titled ``Beauville--Voisin Conjecture for Hyperkähler Manifolds,'' in the Oberwolfach Reports Open Problems collection. No individual proposer is supplied; the proposed year is recorded as 2014.
\item Supplied report citation: \emph{Classical Algebraic Geometry (2015)}, Oberwolfach Reports, DOI 10.4171/owr/2014/31, pp.\ 1695--1745. \url{https://doi.org/10.4171/owr/2014/31}. The supplied background describes Oberwolfach Reports as documenting workshops and their open-problem discussions. The citation's title year 2015 differs from the DOI year and recorded proposed year 2014; this discrepancy has not been resolved.
\item Supplied literature reference: ``Beauville--Voisin--Franchetta conjecture and LLSS eightfolds.'' \url{https://doi.org/10.1016/j.indag.2024.10.004}.
\item Supplied literature reference: ``Voisin, hyperkähler Chow rings.'' \url{https://arxiv.org/abs/math/0602400}.
\item Supplied literature reference: ``Fu, generalized Kummer varieties.'' \url{https://arxiv.org/abs/1309.4977}.
\item Literature assessment dated exactly \texttt{2026-08-21}: the supplied curation labels the problem partially solved and reports proofs for several deformation families and special Hilbert schemes, generalized Kummers, cubic-fourfold Fano varieties, and other cases, while reporting the arbitrary projective hyperkähler case as open. This is supplied dated curation derived from the OWR record, not an independently verified literature assessment; the status and the scope of the cited claims require independent verification.
\end{itemize}
\end{document}
